\akapit Pomiarów dokonywano na specjalnie przygotowanym stanowisku pomiarowym składającym się z następujących elementów:

\begin{itemize}
    \item \textbf{Tunel aerodynaimiczny} - przestrzeń, w której umieszczono zespół napędowy wraz z potrzebnym oprzyrządowaniem, odczyt parametrów tunelu ($q_{tun}$ i $t_{tun}$)
    \item \textbf{Zespół napędowy} - silnik elektryczny + śmigło oraz obwód zasilacza (pozwalający odczytać $U_{zaś}$ i $I_{zaś}$) i zadajnika sygnału PWM
    \item \textbf{Hamownia} - czujniki sił\footnote{belki tensometryczne}, z których odczytu dokonywano w środowisku LabView 
    \item \textbf{Czujniki położenia śmigła} - odczyt obrotów śmigła przez oscyloskop
    \item \textbf{Rurka Prandtla} - odczyt ciśnienia dynamicznego za śmigłem
\end{itemize}
